\section{Scope and Problem Framing}

\subsection{Deliberate Scope}

The project is \textbf{deliberately limited to simulation} in order to increase feasibility and reduce hardware dependence.
This means:
\begin{itemize}
  \item the required deliverable is a working pipeline on standard simulation benchmarks;
  \item real-robot deployment or sim-to-real transfer is not a first-round requirement;
  \item deployment on SO-101 or LeRobot is treated as future work.
\end{itemize}

\subsection{Core Contribution}

The core contribution is:
\begin{quote}
  \textbf{A retrieval-conditioned multimodal reward model that generates dense progress signals for long-horizon robot manipulation in simulation.}
\end{quote}

What is \textbf{not} part of the core contribution:
\begin{itemize}
  \item no training of a foundation model or VLA from scratch;
  \item no requirement for direct continuous action generation from retrieved evidence;
  \item no attempt to prove real-world deployability in the first phase.
\end{itemize}

\section{Proposed Method}

\subsection{Problem Setup}

At each time step $t$, the robot observes the current state $o_t$
(image, state, or both) and the task instruction $g$.
The system also has access to a memory bank $\mathcal{M}$ containing:
\begin{itemize}
  \item instruction or manual text;
  \item stage descriptions;
  \item video clips or sub-trajectory demonstrations;
  \item metadata about the task, object, and outcome.
\end{itemize}

The reward function to be learned is:
\[
r_t = f_\theta(o_t, g, \text{Retrieve}(o_t, g, \mathcal{M}))
\]
where \texttt{Retrieve} returns the most relevant context for the current state.

\subsection{Overall Architecture}

The proposed pipeline has four modules:
\begin{enumerate}
  \item \textbf{Memory construction}: normalize instruction/manual text, stage labels, and segment videos into clips or sub-trajectories.
  \item \textbf{Retriever}: encode the current observation and query top-$k$ relevant evidence.
  \item \textbf{Reward model}: fuse the current observation with retrieved context to estimate progress score or pairwise preference.
  \item \textbf{Policy learner}: use the resulting reward for offline RL, trajectory ranking, or policy selection.
\end{enumerate}

\subsection{Memory Bank Design}

The memory bank should be organized by \textbf{stage} rather than only by full episodes.
A minimal schema for each entry includes:
\begin{itemize}
  \item \texttt{task\_id}, \texttt{stage\_id}, \texttt{instruction}, \texttt{manual\_snippet};
  \item image or clip embeddings;
  \item text embeddings;
  \item outcome labels or relative progress annotations.
\end{itemize}

This design directly supports the key research question on retrieval granularity:
episode-level versus stage-level versus sub-trajectory retrieval.

\subsection{Reward Model}

Two feasible implementation variants are:
\begin{itemize}
  \item \textbf{Progress regression}: predict a continuous or discretized progress score.
  \item \textbf{Pairwise ranking / preference}: given two states or two clips, predict which one is closer to task completion.
\end{itemize}

In the most practical student-scale version of the project, the pretrained encoder remains frozen
or is only lightly fine-tuned, while the main trainable component is the reward head.
This keeps compute manageable and avoids training a large VLM from scratch.

\subsection{How Reward Is Used for Policy Learning}

Three integration levels are possible, in increasing order of ambition:
\begin{enumerate}
  \item use reward first for \textbf{trajectory evaluation} and reward-quality analysis;
  \item use reward for \textbf{offline RL} or \textbf{policy selection};
  \item if resources permit, use reward for post-training a small policy or VLA.
\end{enumerate}

For this proposal, level (2) is the main target because it is strong enough academically while remaining feasible.

\section{Experimental Design}

\subsection{Proposed Benchmarks}

\begin{itemize}
  \item \textbf{Primary benchmark}: \textit{LIBERO} \citep{liu2023libero}, because it is well aligned with long-horizon transfer and multi-task manipulation.
  \item \textbf{Secondary benchmark}: \textit{ManiSkill2/3} \citep{gu2022maniskill2, tao2024maniskill3} when additional evaluation is needed for simulator throughput, task diversity, or scaling behavior.
\end{itemize}

\subsection{Baselines}

\begin{enumerate}
  \item \textbf{Sparse reward only}.
  \item \textbf{Reward model without retrieval}.
  \item \textbf{Retrieval-only guidance} without reward modeling.
  \item \textbf{Retrieval-Augmented Reward Model} (proposed method).
\end{enumerate}

\subsection{Required Ablations}

\begin{itemize}
  \item retrieval source: text only, video only, or multimodal;
  \item granularity: episode, stage, or sub-trajectory;
  \item number of retrieved context items top-$k$;
  \item reward formulation: regression versus ranking;
  \item evaluation split: in-distribution versus OOD.
\end{itemize}

\subsection{Evaluation Metrics}

\begin{itemize}
  \item task-level \textit{success rate};
  \item \textit{normalized return} or cumulative reward;
  \item \textit{sample efficiency};
  \item correlation between predicted reward and ground-truth progress or human preference;
  \item robustness on OOD splits;
  \item retrieval and reward-inference latency.
\end{itemize}

\section{Feasibility Assessment}

\subsection{Technical Feasibility}

If the project tries to combine real robots, OOD detection, retrieval, multimodal reasoning,
action guidance, and hardware deployment in a single phase, the risk becomes too high.
By contrast, \textbf{simulation-first + reward modeling is a realistic scope} because:
\begin{itemize}
  \item the benchmarks and data pipelines already exist;
  \item no foundation model pretraining is required;
  \item retrieval and reward modeling can be debugged as separate modules;
  \item evaluation can be standardized through benchmark suites and ablations.
\end{itemize}

\subsection{Estimated Compute Requirements}

\begin{table}[htbp]
  \centering
  \small
  \begin{tabularx}{\textwidth}{p{2.8cm}p{2.8cm}Y Y}
    \toprule
    Resource tier & Suggested configuration & What it supports & Notes \\
    \midrule
    Minimum & 1 GPU with 12--16GB VRAM, 32GB RAM, 100--200GB SSD & Retrieval prototype, small reward head, subset tasks, and basic ablations & Best with frozen encoders and fewer tasks \\
    Recommended & 1 GPU with 24GB VRAM, 64GB RAM, 200--500GB SSD & Full simulation-only pipeline, more seeds, stage-level retrieval, and offline RL & This is the most appropriate target for the current proposal \\
    Comfortable & A100 40/80GB or equivalent, 64GB+ RAM & Larger sweeps, bigger batch sizes, stronger policy backbones, or extra benchmark coverage & Not required for the core contribution \\
    \bottomrule
  \end{tabularx}
  \caption{Estimated compute requirements by ambition level.}
\end{table}

Under the ``recommended'' configuration, a compact research cycle typically stays within a few GPU-days:
memory-bank preprocessing in a few hours to one day, reward training in a few hours,
policy training in roughly one to three days, and evaluation or ablation in another one to two days.

\subsection{Risks and Mitigation}

\begin{enumerate}
  \item \textbf{Incorrect retrieval context}: use confidence scoring and a fallback to the base reward model when evidence is unreliable.
  \item \textbf{Benchmark overfitting}: include OOD splits and at least one secondary benchmark.
  \item \textbf{Compute overruns}: keep encoders frozen, avoid fine-tuning large VLAs, and prioritize stage-level retrieval first.
  \item \textbf{Tasks are too difficult at full scale}: begin with tasks that have clear stage decomposition, then expand gradually.
\end{enumerate}
